% =====================================================================
%  BlackPearl Morphix -- full research paper (IEEEtran journal format)
%  Compile on Overleaf with pdfLaTeX.
%  FIGURES: upload into the project root, renamed as below (labelled
%  placeholder boxes are shown until then):
%    Stage 1 zip  : real_vs_generated.png, flow_trajectory.png
%    Stage 2 zip  : predicted_vs_observed.png -> s2_predicted_vs_observed.png
%    Step 3       : screenshot of the heatmap   -> s2_signal_heatmap.png
%    Stage 2b zip : r2_by_method.png -> s2b_r2_by_method.png
%                   pred_vs_real.png -> s2b_pred_vs_real.png
%                   examples.png     -> s2b_examples.png
%    Stage 3 zip  : recall_vs_cells.png -> s3_recall_vs_cells.png
%                   design_comparison.png -> s3_design_comparison.png
%  Red text in [brackets] = items still to fill in.
% =====================================================================
\documentclass[journal]{IEEEtran}

\usepackage{amsmath,amssymb}
\usepackage{graphicx}
\usepackage{booktabs}
\usepackage{multirow}
\usepackage{xcolor}
\usepackage{listings}
\usepackage[hidelinks]{hyperref}

\newcommand{\todo}[1]{\textcolor{red}{[#1]}}
\newcommand{\morphixfig}[2]{%
  \IfFileExists{#1}{\includegraphics[width=#2]{#1}}%
  {\fbox{\parbox[c][3.2cm][c]{\dimexpr #2 - 2em\relax}{\centering\small Upload \texttt{\detokenize{#1}}}}}}
\newcommand{\vtheta}{v_\theta}
\newcommand{\E}{\mathbb{E}}

\lstset{language=Python, basicstyle=\ttfamily\scriptsize, numbers=left, numberstyle=\tiny\color{gray},
  breaklines=true, frame=tb, keywordstyle=\color{blue!70!black}, commentstyle=\color{green!40!black},
  stringstyle=\color{red!60!black}, columns=fullflexible, showstringspaces=false}

\markboth{BlackPearl Morphix Research Draft, October 2026}%
{Raj and Srimurugan: Generative Models of Tissue Neighbourhoods}

\begin{document}

\title{BlackPearl Morphix: Generative Models of Tissue Neighbourhoods from Spatial Single-Cell Data, and What They Can and Cannot Learn About Genetic Perturbations}

\author{Aditya Raj and Senthuran Srimurugan%
\thanks{Aditya Raj is with the School of Informatics, The University of Edinburgh, Edinburgh, Scotland, UK.}%
\thanks{Senthuran Srimurugan is with the School of Informatics, The University of Edinburgh, Edinburgh, Scotland, UK.}%
\thanks{Code and results: \todo{GitHub repository URL}.}}

\maketitle

% =====================================================================
\begin{abstract}
We study whether a generative model can learn how cells are arranged in tissue, and what determines that arrangement. BlackPearl Morphix represents a tissue neighbourhood as an unordered set of cells (2D position and cell class) and learns its distribution with conditional flow matching, using a permutation-equivariant transformer with distance-dependent attention. In \emph{Stage~1}, trained on a public mouse-embryo seqFISH dataset and evaluated on spatially held-out tissue, the unconditional model closes $87 \pm 8\%$ of the gap between random scatter and real tissue on cell-type neighbourhood structure (three splits) and outperforms a composition-matched shuffled baseline in every split. In \emph{Stage~2} we use a public Perturb-FISH CRISPR screen in tumour xenografts (187{,}215 cells; 35 knocked-out genes), which we reconstruct into a cell-level table linking position, cell class and knockout. A location-controlled permutation test finds real knockout effects on neighbourhood composition (14 significant tests across 7 knockouts), but neither a knockout-conditioned model nor a simple lookup baseline predicts knockout-specific neighbourhoods in held-out tissue better than assuming no effect: in this pooled design, effects are small, spatially heterogeneous and confounded with location. In contrast, conditioning on each tumour cell's own expression profile, after removing genes that could leak from neighbouring cells, predicts its neighbourhood in held-out tissue (Spearman $\rho = 0.57$ for host-cell fraction and $0.41$ for T-cell fraction; shuffled-expression control $\rho \approx 0$), matching a strong linear baseline and recovering known biology (interferon-response genes near T cells; hypoxia genes away from them). Finally (\emph{Stage~3}), a semi-synthetic virtual laboratory built on the real tumour tissue, with knockout effects planted at the measured sizes, shows why and how such experiments should be designed. Real neighbourhoods are strongly spatially correlated: the mean of 50 cells from one tissue patch varies about five times more than the mean of 50 randomly placed cells, so a mosaic (clonal) layout yields 15 false hits among 22 null knockouts per screen, whereas a randomised layout yields none. With randomised layouts, an active-learning strategy that chooses which knockout to test next finds 80\% of real effects with $2.3\times$ fewer cells than a standard screen (3{,}250 vs.\ 7{,}500) and 90\% of effects in every simulated world within 4{,}500 cells, a level the standard screen reached in only 40 of 100 worlds within 15{,}000 cells. Together, these results argue that learning causal genotype-to-tissue maps requires purpose-designed, randomised, adaptively chosen experiments, which motivates a closed-loop, robot-executed laboratory.
\end{abstract}

\begin{IEEEkeywords}
Flow matching, generative models, spatial transcriptomics, CRISPR screens, Perturb-FISH, tumour microenvironment, permutation equivariance, active learning, experimental design, laboratory automation.
\end{IEEEkeywords}

% =====================================================================
\section{Introduction}

\IEEEPARstart{H}{ow} cells are arranged within a tissue is as informative as what each cell is doing. Immune niches, epithelial sheets and tumour microenvironments all depend on specific cell types sitting next to specific neighbours. Spatial single-cell technologies such as seqFISH and MERFISH now measure, for every cell, a position and a molecular profile~\cite{lohoff2022,palla2022}, and pooled CRISPR screens with imaging readouts add a known genetic perturbation per cell~\cite{binan2025}.

Our long-term question is whether a model can predict how a genetic perturbation changes tissue organisation, so that experiments can be prioritised computationally. This paper takes three steps towards that question and reports each outcome honestly, including a negative one.

\subsection{Contributions}
\begin{itemize}
  \item \textbf{A generative model of tissue neighbourhoods} (Stage 1): cells as an unordered set of positions and types, generated jointly by flow matching~\cite{lipman2023,liu2023} with a permutation-equivariant, distance-aware transformer. On held-out tissue it reproduces cell-type adjacency well beyond composition-matched baselines.
  \item \textbf{A cell-level reconstruction of a public Perturb-FISH screen}~\cite{binan2025} linking every cell's position, class and CRISPR knockout, recovered from separately published tables by exact matching of expression rows.
  \item \textbf{A careful negative result} (Stage 2): knockout effects on neighbourhoods are detectable with location-controlled statistics, but are too small, heterogeneous and location-confounded in this design to be predicted in new tissue by any method we tried, including a simple lookup.
  \item \textbf{A positive result} (Stage 2b): a tumour cell's own expression profile, after a pre-specified leakage filter, predicts its neighbourhood in held-out tissue and recovers known immunological programmes.
  \item \textbf{A simulated closed-loop laboratory} (Stage 3) on real tissue, showing that spatial correlation makes mosaic layouts unreliable and that active experiment selection with randomised layouts finds real effects with several-fold fewer cells, providing a quantitative design for robot-executed perturbation studies.
\end{itemize}

\subsection{Scope}
All models work on 2D neighbourhoods of 64 cells. This is a research tool for hypothesis generation; it makes no clinical or diagnostic claims and does not replace experimental validation.

% =====================================================================
\section{Related Work}

\textbf{Spatial mapping of single-cell data.} Tangram~\cite{biancalani2021} and novoSpaRc~\cite{nitzan2019} place dissociated single-cell profiles onto spatial coordinates; they map existing cells rather than generating new arrangements. Graph-based models relate a cell's expression to its spatial neighbours~\cite{fischer2023}.

\textbf{Generative models of sets.} Diffusion~\cite{ho2020} and flow-matching~\cite{lipman2023,liu2023,albergo2023,tong2024} models transport noise to data. Set architectures~\cite{zaheer2017,lee2019} and equivariant networks~\cite{satorras2021,hoogeboom2022} make generation invariant to element order, as in molecule generation where atoms are positioned and typed jointly~\cite{hoogeboom2022}. Our setting is analogous with cells in place of atoms.

\textbf{Perturbation modelling.} MorphoDiff~\cite{navidi2025} and CellFlux~\cite{zhang2025} generate images of single cells under perturbations, trained on Cell Painting resources~\cite{bray2016,chandrasekaran2023}. Perturb-FISH~\cite{binan2025} combines pooled CRISPR screening with MERFISH imaging, including in tumours. We model multi-cell spatial arrangement rather than single-cell appearance.

\textbf{Self-driving laboratories.} Closed loops in which a model chooses experiments and a robot executes them have been demonstrated for functional genomics~\cite{king2004} and chemistry~\cite{burger2020}, often guided by active learning~\cite{settles2009}.

% =====================================================================
\section{Method}

\subsection{Representation}
A neighbourhood is a set of $N$ cells $\{(\mathbf{p}_i, c_i)\}$ with positions $\mathbf{p}_i\in\mathbb{R}^2$ and classes $c_i\in\{1,\dots,K\}$, encoded as $\mathbf{x}^{(i)} = [\mathbf{p}_i/\sigma,\ \mathrm{onehot}(c_i)]$, where $\sigma$ is the standard deviation of training positions. Positions are measured in \emph{cell spacings} (coordinates rescaled so that the median nearest-neighbour distance is 1). Random rotations and reflections are used as augmentation.

\subsection{Conditional Flow Matching}
A velocity field $\vtheta(\mathbf{x}_t, t \,|\, \mathbf{z})$ transports noise $\mathbf{x}_0\sim\mathcal{N}(\mathbf{0},\mathbf{I})$ to data $\mathbf{x}_1$ along straight paths $\mathbf{x}_t=(1-t)\mathbf{x}_0+t\mathbf{x}_1$~\cite{lipman2023,liu2023}, trained with
\begin{equation}
\begin{split}
\mathcal{L}_{\mathrm{FM}} = {}& \E\big\|\vtheta^{\mathrm{pos}} - (\mathbf{x}_1-\mathbf{x}_0)^{\mathrm{pos}}\big\|^2 \\
&+ \E\big\|\vtheta^{\mathrm{type}} - (\mathbf{x}_1-\mathbf{x}_0)^{\mathrm{type}}\big\|^2.
\end{split}
\end{equation}
Positions and classes are generated jointly; each class is read out as the argmax of its channels. Samples are drawn by Euler integration (50--100 steps). The condition $\mathbf{z}$ is empty in Stage~1, a knockout label in Stage~2 and an expression profile in Stage~2b.

\subsection{Network}
Each cell is embedded linearly and processed by 4 transformer blocks (width 128, 4 heads) with no index positional encoding, so the network is permutation-equivariant. Time (and the condition) modulate every block through adaptive layer normalisation~\cite{peebles2023}. Spatial structure enters as an additive attention bias computed from current pairwise distances, $\beta^{(h)}_{ij} = \mathbf{w}_h^{\top}\boldsymbol{\phi}(\|\mathbf{p}_i-\mathbf{p}_j\|)$ with 16 radial basis functions, so each block is message passing on a distance-weighted cell graph (1.2\,M parameters).

\subsection{Conditioning (Stages 2 and 2b)}
Neighbourhoods are centred on a \emph{focal} tumour cell, which enters the network as a fixed token at the origin so that generated neighbours can respect its position. The condition is either a learned knockout embedding (Stage~2) or a two-layer projection of the focal cell's standardised log-expression (Stage~2b), added to the time embedding. During training the condition is replaced by a learned ``null'' embedding with probability 0.10--0.15, so one network is both the conditional and the unconditional model.

\subsection{Physics Prior (Stage 1 Ablation)}
An optional penalty on the predicted endpoint $\hat{\mathbf{p}}_1 = \mathbf{p}_t + (1-t)\vtheta^{\mathrm{pos}}$ discourages cells closer than $d_{\min}$ (the 2nd percentile of real nearest-neighbour distances):
\begin{equation}
\mathcal{L}_{\mathrm{ov}} = \E_t\Big[t\cdot\tfrac{1}{N}\textstyle\sum_{i\neq j}\max\big(0,\,d_{\min}-\|\hat{\mathbf{p}}_1^{(i)}-\hat{\mathbf{p}}_1^{(j)}\|\big)^2\Big].
\end{equation}

\subsection{Training}
AdamW~\cite{loshchilov2019} (learning rate $3\times10^{-4}$, weight decay 0.01), 500 warm-up steps and cosine decay, batch 128, 12{,}000--15{,}000 steps, gradient clipping 1.0, and an exponential moving average of weights (0.999) for evaluation. Each model trains in about 10 minutes on one NVIDIA T4 GPU.

\subsection{Evaluation Principles}
All evaluations use \textbf{spatially disjoint} train/test splits: tissue is cut into a grid of blocks and whole blocks are held out, so test neighbourhoods never share cells with training neighbourhoods. Every result is compared with simple baselines, and in Stages~2 and~2b the test, controls and success criteria were fixed before the runs.

% =====================================================================
\section{Stage 1: Learning Tissue Organisation}

\subsection{Data and Protocol}
We use the mouse-embryo seqFISH dataset of Lohoff \emph{et al.}~\cite{lohoff2022} via Squidpy~\cite{palla2022}: 19{,}416 cells with 22 annotated cell types. The tissue is split into a $4\times4$ grid and blocks holding $\approx$20\% of cells are held out. A patch is a cell and its 63 nearest neighbours, kept only if all 64 cells lie in the same split (20{,}000 training and 2{,}000 test patches). Each of three seeds changes both the model initialisation and the held-out blocks.

Generated patches are compared with 500 held-out patches using Wasserstein-1 distances of nearest-neighbour distance, pairwise distance and patch radius; total variation distance (TVD) of cell-type composition; \textbf{neighbourhood TVD}, the frequency-weighted TVD between real and generated distributions $P(b\,|\,a)$ of a type-$a$ cell's 5 nearest neighbours having type $b$; and the overlap rate. References: \emph{random scatter} (uniform disc, independent types); \emph{type-shuffled real} (real geometry and composition, types permuted within each patch); and \emph{real train vs.\ test} (another real region; not an upper bound, because embryo anatomy is region-specific).

\subsection{Results}

\begin{table*}[!t]
\centering
\caption{Stage 1: generated vs.\ held-out real tissue, mean $\pm$ s.d.\ over three seeds (each a different held-out split). Lower is better. The last row counts seeds in which the model beat random scatter.}
\label{tab:s1}
\small
\begin{tabular}{lcccccc}
\toprule
Condition & NN-dist.\ $W_1$ & Pair-dist.\ $W_1$ & Radius $W_1$ & Comp.\ TVD & Neigh.\ TVD & Overlap \\
\midrule
Random scatter          & $0.120 \pm 0.087$ & $0.577 \pm 0.132$ & $0.821 \pm 0.033$ & $0.406 \pm 0.088$ & $0.652 \pm 0.035$ & $0.022 \pm 0.001$ \\
Type-shuffled real      & $0.076 \pm 0.071$ & $0.423 \pm 0.358$ & $0.370 \pm 0.270$ & $0.406 \pm 0.088$ & $0.328 \pm 0.056$ & $0.020 \pm 0.002$ \\
Real train vs.\ test    & $0.078 \pm 0.068$ & $0.421 \pm 0.224$ & $0.381 \pm 0.258$ & $0.401 \pm 0.094$ & $0.238 \pm 0.053$ & $0.020 \pm 0.000$ \\
\midrule
\textbf{Morphix} & $0.090 \pm 0.078$ & $0.450 \pm 0.307$ & $0.385 \pm 0.254$ & $0.390 \pm 0.079$ & $\mathbf{0.293 \pm 0.070}$ & $0.022 \pm 0.002$ \\
\quad beats random (seeds) & 3/3 & 2/3 & 3/3 & 3/3 & 3/3 & -- \\
\bottomrule
\end{tabular}
\end{table*}

Because each seed holds out a different region, we rely on paired comparisons within a split (Table~\ref{tab:s1}). On neighbourhood structure the model closes $87 \pm 8\%$ of the random-to-real gap (93\%, 90\%, 77\%) and beats the type-shuffled baseline in all three seeds (by 0.026, 0.053, 0.024), recovering $41 \pm 19\%$ of the remaining shuffled-to-real gap. It therefore learns which cell types are adjacent, not only the overall mix. The weakest split (seed~2) is the one whose held-out blocks differ most from training in cell-type mix (composition TVD 0.50). The overlap rate matches real tissue. Geometric metrics vary more between splits than between methods, and two disjoint sets of 500 real patches can differ by 0.18 in pair-distance $W_1$, so we draw no fine conclusions from them; generated patches appear slightly more uniformly packed than real ones (Fig.~\ref{fig:s1}).

An explicit overlap penalty gave no measurable benefit, even when training on only 500 patches (overlap rate 0.024 with vs.\ 0.023 without; the penalty stayed near $10^{-4}$): plausible spacing is learned from data alone. Reducing training data from 20{,}000 to 500 patches worsened neighbourhood TVD from 0.298 to $\approx$0.36 (confounded with fewer training steps).

\begin{figure*}[!t]
\centering
\morphixfig{real_vs_generated.png}{0.95\textwidth}
\caption{Stage 1. Top: real held-out patches. Bottom: patches generated from noise. Colours are cell types; columns are not paired.}
\label{fig:s1}
\end{figure*}

% =====================================================================
\section{Stage 2: A Spatial CRISPR Screen}

\subsection{Data Reconstruction}
Perturb-FISH~\cite{binan2025} imaged human A375 melanoma cells carrying pooled CRISPR knockouts of genes in the TLR4/NF-$\kappa$B pathway, grown as xenograft tumours in mice together with human T cells, using MERFISH with a 500-gene panel. The processed tables (Brain Image Library) store positions and guide calls without gene names, and gene-named perturbation tables without positions. We reconstructed a cell-level table by (i) reading the header-less position and guide tables (187{,}215 cells each), (ii) locating the 500 gene columns within a 553-column count table (the remaining columns hold metadata and blank barcodes) using anchor cells whose count rows occur exactly once, and (iii) matching each perturbed cell's count row exactly to the full table. All 9{,}432 perturbed tumour cells and 14{,}232 T cells matched uniquely (100\%). As an independent check, each of the 77 unnamed guide columns mapped almost exclusively to a single gene (typically two guides per gene).

Tumour cells carry 35 knockouts plus a non-targeting Control group (4{,}184 cells); knockout group sizes range from 2 to 416 cells. T-cell knockouts are too sparse to model (12{,}966 of 14{,}232 are Control), so we study \textbf{tumour-cell knockouts} and the cells around them. Same-knockout cells are moderately clustered: 39\% of a perturbed tumour cell's 5 nearest perturbed neighbours share its knockout, vs.\ 23\% expected at random. Neighbour classes are \emph{tumour-like} (perturbed tumour cells and unlabelled cells with tumour-like total counts), \emph{T cell}, and \emph{host} (unlabelled cells with low human-probe counts, mostly mouse; threshold = 5th percentile of tumour-cell log counts).

\subsection{Do Knockouts Change Their Neighbourhood?}
For 8{,}911 single-knockout and Control tumour cells we computed six features of the 63 nearest neighbours (fractions of T-cell, host and tumour-like neighbours; mean and nearest distance; distance to the nearest T cell). Each knockout with $\geq$30 cells (27 knockouts, 162 tests) was compared with Control cells \textbf{in the same tissue block} ($\approx$40 cell spacings) using 2{,}000 block-stratified label permutations, with Benjamini--Hochberg correction. Stratification prevents a knockout from appearing T-cell rich merely because its cells sit in T-cell-rich regions; on synthetic data with a planted effect and a planted location confound, the test detected the former ($q=0.015$) and correctly ignored the latter ($p=0.65$).

\textbf{Result.} 14 tests across 7 knockouts were significant ($q<0.1$): IRF5, IRF7, LBP, MAP2K3, MAP2K6, NFKBIA and PELI1 (Fig.~\ref{fig:signal}). The p-value distribution is enriched near zero (33 of 162 below 0.1, vs.\ $\approx$16 expected under no effect). The largest effect is NFKBIA: knocked-out cells lie 2.7 cell spacings further from the nearest T cell than nearby Control cells (1.05 SD) and have fewer T-cell neighbours. Several knockouts (IRF5, LBP, MAP2K3, MAP2K6) shift towards more host and fewer tumour-like neighbours. Most effects are modest (0.2--0.6 SD).

\begin{figure}[!t]
\centering
\morphixfig{s2_signal_heatmap.png}{\columnwidth}
\caption{Knockout vs.\ same-block Control neighbourhoods (effect size in SD units; $*$: $q<0.1$).}
\label{fig:signal}
\end{figure}

\subsection{Can a Model Predict Knockout-Specific Neighbourhoods?}
We trained the conditional model on 63-cell neighbourhoods around 6{,}423 labelled training tumour cells (28 conditions) plus 40{,}000 unlabelled tumour-like cells under the null condition, holding out $6\times6$ tissue blocks ($\approx$20\%; 1{,}811 labelled test cells). We evaluated (i) the held-out flow-matching loss on knockout cells given their true label vs.\ a Control label, and (ii) for each knockout and feature, the predicted shift relative to Control vs.\ the shift observed in held-out tissue, both raw and \textbf{location-matched} (knockout vs.\ Control cells in the same block). Baselines: \emph{no effect} (zero shift) and \emph{lookup} (the shift measured in training tissue). The \emph{noise floor} is the error a perfect predictor would make due to sampling error in the small held-out groups.

\begin{table}[!t]
\centering
\caption{Stage 2: predicting knockout effects in held-out tissue (21--22 knockouts $\times$ 5 features; mean squared error in Control-SD units; $r$: Pearson correlation with observed shifts).}
\label{tab:s2}
\footnotesize
\setlength{\tabcolsep}{4pt}
\begin{tabular}{lcc|cc}
\toprule
 & \multicolumn{2}{c|}{Raw shifts} & \multicolumn{2}{c}{Location-matched} \\
 & MSE & $r$ & MSE & $r$ \\
\midrule
Noise floor          & 0.068 & -- & 0.067 & -- \\
No effect (zero)     & 0.362 & -- & \textbf{0.096} & -- \\
Lookup (train shift) & 0.362 & 0.18 & 0.106 & 0.05 \\
Morphix              & 0.349 & 0.21 & 0.129 & 0.06 \\
\bottomrule
\end{tabular}
\end{table}

\textbf{Result: no generalisable knockout signal.} The held-out loss on knockout cells was the same with the true label (1.7874), a Control label (1.7865) or no label (1.7866). After location matching, assuming \emph{no effect} predicted held-out shifts best (Table~\ref{tab:s2}, Fig.~\ref{fig:s2}); both the model and the simple lookup did worse, and the zero-effect error was close to the noise floor. Raw (unmatched) shifts were large but did not replicate between regions, consistent with location confounding. The limit is therefore the data rather than the model: the lookup baseline, which uses the training effects directly, also fails. NFKBIA is a partial exception: in held-out tissue its location-matched effects were strong (distance to T cell $+0.88 \pm 0.23$ SD; T-cell fraction $-0.69 \pm 0.11$ SD) and in the same direction as the whole-tissue test, but weak in the training blocks ($+0.14$, $+0.09$), indicating a real but spatially heterogeneous effect.

We attribute the null result to three properties of the design: few cells per knockout (2--416 in total; roughly 20--80 per knockout in a held-out region), a mosaic layout in which each cell's neighbours carry different knockouts, and clonal clustering that ties knockout to location. Detecting a 0.2 SD effect at about two standard errors needs a standard error of $\approx$0.1 SD, i.e.\ roughly 200 knockout cells compared with a similar number of location-matched controls, which few knockouts here reach.

\begin{figure}[!t]
\centering
\morphixfig{s2_predicted_vs_observed.png}{\columnwidth}
\caption{Stage 2, location-matched: predicted vs.\ observed knockout shifts in held-out tissue for Morphix (left) and the lookup baseline (right). Error bars: held-out standard errors.}
\label{fig:s2}
\end{figure}

% =====================================================================
\section{Stage 2b: Cell State Predicts Neighbourhood}

\subsection{Question and Pre-Specified Test}
Does a tumour cell's own expression profile predict the neighbourhood it sits in? Fixed before running: on held-out blocks, for 1{,}500 test cells, the model generates 16 neighbourhoods from the cell's expression; their mean features are compared with the cell's real neighbourhood (R$^2$, Spearman $\rho$). Success required (1) R$^2>0$ for T-cell and host fractions, (2) a \textbf{shuffled-expression} control near zero, and (3) lower held-out loss with the cell's own expression than shuffled. Baselines: no expression (training mean), $k$-nearest neighbours ($k=50$) in a 32-dimensional PCA of expression, and ridge regression on all kept genes.

\textbf{Leakage guard.} In imaging data a cell touching a T cell can acquire some of that cell's transcripts through imperfect segmentation, which would let a model ``detect'' neighbours trivially. We removed every gene whose mean log-normalised expression in T cells or host cells exceeded twice that in tumour cells: 56 of 500 genes, including T-cell markers (TRAC, CD3D/E/G, CD8A, CD2, GZMB, PRF1, PTPRC) and stromal genes (COL1A1, PECAM1). A second model using all 500 genes measures how much leakage could inflate results.

The model was trained on 60{,}000 tumour-like focal cells (15{,}000 test) with the $6\times6$ block split; the condition is a projection of the standardised log-expression of the kept genes.

\subsection{Results}

\begin{table}[!t]
\centering
\caption{Stage 2b: predicting each held-out cell's real neighbourhood from its expression (clean genes; one split). R$^2$ / Spearman $\rho$.}
\label{tab:s2b}
\footnotesize
\setlength{\tabcolsep}{3pt}
\begin{tabular}{lccc}
\toprule
Method & T-cell frac. & Host frac. & Dist.\ to T cell \\
\midrule
No expression       & 0.00 / -- & 0.00 / -- & 0.00 / -- \\
Shuffled expression & $-0.33$ / 0.02 & $-0.60$ / $-0.04$ & $-0.20$ / 0.01 \\
$k$NN ($k=50$)      & 0.18 / 0.39 & 0.20 / 0.55 & 0.15 / 0.37 \\
Ridge regression    & 0.17 / 0.41 & 0.40 / 0.57 & 0.15 / 0.38 \\
\textbf{Morphix}    & 0.16 / 0.41 & 0.40 / 0.57 & 0.13 / 0.36 \\
\bottomrule
\end{tabular}
\end{table}

All three pre-specified criteria were met (Table~\ref{tab:s2b}, Fig.~\ref{fig:s2b}). Expression predicts host fraction (R$^2=0.40$, $\rho=0.57$), T-cell fraction (R$^2=0.16$, $\rho=0.41$) and distance to the nearest T cell ($\rho=0.36$) in unseen tissue. The shuffled-expression control has $\rho\approx0$ (its negative R$^2$ reflects confident predictions assigned to the wrong cells), and the held-out loss was lowest with the cell's own expression (1.716) vs.\ none (1.721) and shuffled (1.727). Using all genes changed little (T-cell fraction R$^2$ 0.17 vs.\ 0.16; host 0.38 vs.\ 0.40), so leaking transcripts do not drive the result.

Morphix matches ridge regression on composition and outperforms $k$NN on host fraction, but is worse on spacing features (R$^2$ $-0.12$ and $-0.56$ for mean and nearest distance, vs.\ 0.08 and 0.14 for ridge; ranking is preserved, $\rho\approx0.25$), reflecting a bias in generated distances also seen in Stage~2. Its distinct contribution is generating complete, cell-resolved neighbourhoods rather than summary predictions (Fig.~\ref{fig:s2bex}). These results come from one split \todo{update with mean $\pm$ s.d.\ over seeds 0--2}.

\textbf{Biological consistency (associations, not causes).} Ridge coefficients on the kept genes associate T-cell-rich neighbourhoods with CXCL10, CXCL11, IDO1, HLA-B and TAPBP, consistent with the interferon-$\gamma$ response of tumour cells to nearby T cells~\cite{tokunaga2018}, and T-cell-poor neighbourhoods with PKM, VEGFA, PDK1 and CA9, consistent with hypoxic, glycolytic regions from which T cells are often excluded \todo{citation}. None of these genes was removed by the leakage filter, so they reflect the tumour cells' own state. Because neighbourhoods also shape expression, these links are associative.

\begin{figure}[!t]
\centering
\morphixfig{s2b_r2_by_method.png}{\columnwidth}
\caption{Stage 2b: R$^2$ on held-out cells by method and feature, for clean genes (left) and all genes (right).}
\label{fig:s2b}
\end{figure}

\begin{figure*}[!t]
\centering
\morphixfig{s2b_examples.png}{0.95\textwidth}
\caption{Stage 2b: held-out cells (gold star) with the highest (left) and lowest (right) predicted T-cell fraction. Top: real neighbourhood. Bottom: one neighbourhood generated from the cell's expression. Red: tumour-like; blue: T cell; grey: host.}
\label{fig:s2bex}
\end{figure*}

% =====================================================================
\section{Stage 3: Simulated Closed-Loop Experiment Selection}
\label{sec:stage3}

Stage~2 suggests that the bottleneck is experimental design rather than modelling. We therefore asked, in a setting where the truth is known: \emph{which physical layout and which experiment-selection strategy recover knockout effects most efficiently?}

\subsection{Virtual Laboratory}
The virtual laboratory draws every measured cell from \textbf{real} tissue: 40{,}000 tumour-like cells from the Perturb-FISH sections, each described by three standardised neighbourhood features (T-cell fraction, host fraction, distance to nearest T cell; same definitions as Stage~2). Cell-to-cell variability and spatial structure are therefore real. On top of these, each simulated ``world'' plants knockout effects at the sizes measured in Stage~2: of 30 knockouts, 8 shift one randomly chosen feature by $\pm$0.2, 0.25, 0.3, 0.35, 0.4, 0.5, 0.6 or 1.0~SD, and 22 have no effect. One \emph{experiment} measures 50 cells carrying one knockout; the budget is 300 experiments (15{,}000 cells). Two physical layouts are compared: \emph{randomised} (cells drawn from random positions, as in an arrayed, randomised design) and \emph{mosaic} (the 50 cells nearest a random point, as for a clone in a pooled tumour screen).

\subsection{Inference and Strategies}
Each knockout--feature effect $\theta$ has a Gaussian posterior under a $\mathcal{N}(0, 0.5^2)$ prior and unit cell-level noise: after $n$ cells with summed readout $s$, the posterior mean and variance are $m = s/(0.5^{-2}+n)$ and $v = 1/(0.5^{-2}+n)$. A knockout is called a \textbf{hit} when $P(|\theta|>0.1\,\mathrm{SD}) > 0.99$ for some feature \emph{and} it has been measured in at least 150 cells (confirmation by repeats). Bayesian calls remain valid when later experiments depend on earlier results. All strategies use the same rule, and each starts with one screening experiment per knockout:
\begin{itemize}
  \item \textbf{Round robin:} test every knockout in turn (a standard screen).
  \item \textbf{Random:} choose knockouts at random.
  \item \textbf{Active:} test the not-yet-confirmed knockout with the highest posterior probability of a real effect, skipping knockouts that already look confidently null ($P<0.05$).
\end{itemize}
The active strategy and calling rule were developed on synthetic mock tissue \emph{before} running on the real data. A first, uncertainty-sampling strategy (testing the knockout whose hit status was most uncertain) did not beat round robin, because it spent experiments confirming nulls; we report it as a negative design result. Results are averaged over 100 worlds (50 for the harder and design comparisons).

\subsection{Results}

\begin{table}[!t]
\centering
\caption{Stage 3: experiment-selection strategies on real tissue with planted effects (randomised layout; 100 worlds; budget 15{,}000 cells).}
\label{tab:s3}
\footnotesize
\setlength{\tabcolsep}{3pt}
\begin{tabular}{lccc}
\toprule
 & Active & Round robin & Random \\
\midrule
Cells to 80\% recall (median) & \textbf{3{,}250} & 7{,}500 & 9{,}250 \\
Cells to 90\% recall (median) & \textbf{4{,}500} & $>$15{,}000 & $>$15{,}000 \\
Worlds reaching 90\% & \textbf{100/100} & 40/100 & 37/100 \\
Recall at 15{,}000 cells & \textbf{1.00} & 0.91 & 0.89 \\
False hits per screen & 0.03 & 0.00 & 0.00 \\
Effect-size error (SD) & 0.09 & 0.06 & 0.06 \\
\midrule
\multicolumn{4}{l}{\emph{Effects halved:}} \\
Recall at 15{,}000 cells & \textbf{0.81} & 0.44 & -- \\
Cells to 50\% recall & \textbf{3{,}375} & 11{,}625 & -- \\
\bottomrule
\end{tabular}
\end{table}

\textbf{Spatial correlation makes mosaic layouts unreliable.} Without any planted effect, the mean of 50 randomly placed cells varies by 0.14~SD, exactly as expected for independent cells ($1/\sqrt{50}=0.141$). The mean of 50 cells from one contiguous patch varies by 0.59--0.73~SD, about five times more. Consequently, the same round-robin screen and analysis run on a mosaic layout called 15.3 of the 22 null knockouts as hits per screen (false discovery rate 68\%), against none with a randomised layout, while recall was similar (0.90 vs.\ 0.92; Fig.~\ref{fig:s3design}). This quantifies, with real tissue, the location confounding that defeated Stage~2. An analysis that models patch-to-patch variation would reduce these false hits, but at a large cost in power; randomised placement avoids the problem.

\textbf{Active selection finds effects with fewer cells.} With randomised layouts, the active strategy found 80\% of real effects with $2.3\times$ fewer cells than round robin, and reached 90\% in every world within 4{,}500 cells, whereas round robin reached 90\% in only 40 of 100 worlds within 15{,}000 cells (Table~\ref{tab:s3}, Fig.~\ref{fig:s3}). The advantage grew when effects were halved: active selection found 81\% of effects at the full budget against 44\% for round robin, and reached 50\% with $3.4\times$ fewer cells. The costs were small: about 0.03 false hits per screen and slightly less precise effect estimates (0.09 vs.\ 0.06~SD), because active selection stops measuring a knockout once it is confirmed.

\begin{figure}[!t]
\centering
\morphixfig{s3_recall_vs_cells.png}{\columnwidth}
\caption{Stage 3: fraction of planted effects found (left; band: 10--90\% of worlds) and false hits (right) as a function of cells measured, for active, round-robin and random selection.}
\label{fig:s3}
\end{figure}

\begin{figure}[!t]
\centering
\morphixfig{s3_design_comparison.png}{\columnwidth}
\caption{Stage 3: the same screen and analysis with a mosaic (clonal) vs.\ randomised layout. Left: real effects found. Right: false hits among 22 null knockouts.}
\label{fig:s3design}
\end{figure}

% =====================================================================
\section{Toward Designed Experiments and a Closed Robotic Loop}
\label{sec:robotics}

Stages 2--3 point to the same conclusion. Tissue organisation is strongly linked to cell state and is learnable from tens of thousands of cells, but the effect of a single gene knockout is small and, in a pooled mosaic screen, entangled with location and spread over too few cells to learn. Stage~3 shows that the remedy is experimental: randomised layouts remove the location confound, and adaptive selection of which knockout to test next recovers effects with several-fold fewer cells.

Laboratory automation can generate such data, and a generative model can decide which data is most informative. We propose a closed loop in the spirit of self-driving laboratories~\cite{king2004,burger2020}:
\begin{enumerate}
  \item \textbf{Predict:} the conditional model generates neighbourhoods for candidate perturbations.
  \item \textbf{Select:} the next perturbations are chosen by active learning~\cite{settles2009}, e.g.\ the most promising not-yet-confirmed candidates, as validated in Stage~3.
  \item \textbf{Execute:} a collaborative robot arm (e.g.\ a Universal Robots UR3e or UR5e) moves plates between a liquid handler, an incubator and an automated microscope, applying arrayed perturbations (one perturbation per well, in a randomised layout) and imaging the result.
  \item \textbf{Learn:} images are segmented into cell positions and classes, the representation the model already uses, and added to training.
\end{enumerate}
Automation provides scale, consistent handling, randomised layouts and the tight coupling between model decisions and experiments. Stage~3 validates the selection logic in simulation; a physical implementation would require an established laboratory providing the necessary facilities, biosafety approval and supervision, and would first need to confirm that real knockout effects in arrayed cultures resemble the planted additive shifts assumed here.

% =====================================================================
\section{Limitations}
\begin{itemize}
  \item 2D, 64-cell neighbourhoods; no 3D structure or whole tissues.
  \item One embryo section (Stage~1) and one xenograft experiment (Stages~2--2b); held-out regions differ anatomically from training regions, so ``real vs.\ real'' is not a ceiling.
  \item Neighbour classes in Stages~2--2b are coarse (tumour-like, T cell, host), and the host class is a count-based heuristic.
  \item Stage~2b is one split \todo{update}; Stage~2 is one split; Stage~1 uses three.
  \item Generated distances are biased in Stages~2--2b; cell classes use a continuous relaxation and argmax.
  \item Stage~2b associations are not causal, and the leakage filter is a heuristic that cannot exclude all segmentation artefacts.
  \item Stage~3 is semi-synthetic: cells and their variability are real, but knockout effects are planted additive shifts on three summary features with assumed sizes and prevalence; it validates experiment-selection logic, not biology. The active strategy was developed on mock tissue, and its calling rule assumes independent cells, which holds for randomised but not mosaic layouts.
\end{itemize}

% =====================================================================
\section{Conclusion}
A permutation-equivariant flow-matching model learns realistic tissue neighbourhoods and their cell-type adjacency from spatial single-cell data. Conditioned on a cell's own expression, it predicts that cell's neighbourhood in unseen tissue and recovers known immunological programmes. Conditioned on CRISPR knockouts from a pooled spatial screen, neither it nor a simple lookup predicts knockout-specific neighbourhoods beyond what location explains, even though a location-controlled test detects real effects. A semi-synthetic virtual laboratory on real tissue explains why: neighbourhoods are spatially correlated, so mosaic layouts generate many false hits, whereas randomised layouts combined with active experiment selection find real effects with several-fold fewer cells. We conclude that learning how genetic perturbations reshape tissue requires purpose-designed, randomised and adaptively chosen experiments, and that a closed-loop, robot-executed laboratory is a natural way to generate them.

\section*{Acknowledgment}
We thank the authors of the seqFISH and Perturb-FISH studies for making their data public. Computation used Google Colab.

% =====================================================================
\begin{thebibliography}{99}
\bibitem{lohoff2022} T.~Lohoff, S.~Ghazanfar, A.~Missarova \emph{et al.}, ``Integration of spatial and single-cell transcriptomic data elucidates mouse organogenesis,'' \emph{Nature Biotechnology}, vol.~40, pp.~74--85, 2022.
\bibitem{palla2022} G.~Palla, H.~Spitzer, M.~Klein \emph{et al.}, ``Squidpy: a scalable framework for spatial omics analysis,'' \emph{Nature Methods}, vol.~19, pp.~171--178, 2022.
\bibitem{binan2025} L.~Binan \emph{et al.}, ``Simultaneous CRISPR screening and spatial transcriptomics reveal intracellular, intercellular, and functional transcriptional circuits,'' \emph{Cell}, 2025. \todo{verify authors, volume, pages}
\bibitem{lipman2023} Y.~Lipman, R.~T.~Q.~Chen, H.~Ben-Hamu, M.~Nickel, and M.~Le, ``Flow matching for generative modeling,'' in \emph{Proc.\ ICLR}, 2023.
\bibitem{liu2023} X.~Liu, C.~Gong, and Q.~Liu, ``Flow straight and fast: Learning to generate and transfer data with rectified flow,'' in \emph{Proc.\ ICLR}, 2023.
\bibitem{albergo2023} M.~S.~Albergo and E.~Vanden-Eijnden, ``Building normalizing flows with stochastic interpolants,'' in \emph{Proc.\ ICLR}, 2023.
\bibitem{tong2024} A.~Tong, K.~Fatras, N.~Malkin \emph{et al.}, ``Improving and generalizing flow-based generative models with minibatch optimal transport,'' \emph{Transactions on Machine Learning Research}, 2024.
\bibitem{peebles2023} W.~Peebles and S.~Xie, ``Scalable diffusion models with transformers,'' in \emph{Proc.\ ICCV}, 2023.
\bibitem{ho2020} J.~Ho, A.~Jain, and P.~Abbeel, ``Denoising diffusion probabilistic models,'' in \emph{Proc.\ NeurIPS}, 2020.
\bibitem{zaheer2017} M.~Zaheer, S.~Kottur, S.~Ravanbakhsh \emph{et al.}, ``Deep sets,'' in \emph{Proc.\ NeurIPS}, 2017.
\bibitem{lee2019} J.~Lee, Y.~Lee, J.~Kim \emph{et al.}, ``Set transformer: A framework for attention-based permutation-invariant neural networks,'' in \emph{Proc.\ ICML}, 2019.
\bibitem{satorras2021} V.~G.~Satorras, E.~Hoogeboom, and M.~Welling, ``E($n$) equivariant graph neural networks,'' in \emph{Proc.\ ICML}, 2021.
\bibitem{hoogeboom2022} E.~Hoogeboom, V.~G.~Satorras, C.~Vignac, and M.~Welling, ``Equivariant diffusion for molecule generation in 3D,'' in \emph{Proc.\ ICML}, 2022.
\bibitem{biancalani2021} T.~Biancalani, G.~Scalia, L.~Buffoni \emph{et al.}, ``Deep learning and alignment of spatially resolved single-cell transcriptomes with Tangram,'' \emph{Nature Methods}, vol.~18, pp.~1352--1362, 2021.
\bibitem{nitzan2019} M.~Nitzan, N.~Karaiskos, N.~Friedman, and N.~Rajewsky, ``Gene expression cartography,'' \emph{Nature}, vol.~576, pp.~132--137, 2019.
\bibitem{fischer2023} D.~S.~Fischer, A.~C.~Schaar, and F.~J.~Theis, ``Modeling intercellular communication in tissues using spatial graphs of cells,'' \emph{Nature Biotechnology}, vol.~41, pp.~332--336, 2023. \todo{verify}
\bibitem{navidi2025} Z.~Navidi, J.~Ma, E.~A.~Miglietta \emph{et al.}, ``MorphoDiff: Cellular morphology painting with diffusion models,'' in \emph{Proc.\ ICLR}, 2025.
\bibitem{zhang2025} Y.~Zhang, Y.~Su, C.~Wang \emph{et al.}, ``CellFlux: Simulating cellular morphology changes via flow matching,'' in \emph{Proc.\ ICML}, 2025.
\bibitem{bray2016} M.-A.~Bray, S.~Singh, H.~Han \emph{et al.}, ``Cell Painting, a high-content image-based assay for morphological profiling using multiplexed fluorescent dyes,'' \emph{Nature Protocols}, vol.~11, pp.~1757--1774, 2016.
\bibitem{chandrasekaran2023} S.~N.~Chandrasekaran, J.~Ackerman, E.~Alix \emph{et al.}, ``JUMP Cell Painting dataset: morphological impact of 136,000 chemical and genetic perturbations,'' \emph{bioRxiv}, 2023.
\bibitem{tokunaga2018} R.~Tokunaga, W.~Zhang, M.~Naseem \emph{et al.}, ``CXCL9, CXCL10, CXCL11/CXCR3 axis for immune activation -- a target for novel cancer therapy,'' \emph{Cancer Treatment Reviews}, vol.~63, pp.~40--47, 2018. \todo{verify}
\bibitem{king2004} R.~D.~King, K.~E.~Whelan, F.~M.~Jones \emph{et al.}, ``Functional genomic hypothesis generation and experimentation by a robot scientist,'' \emph{Nature}, vol.~427, pp.~247--252, 2004.
\bibitem{burger2020} B.~Burger, P.~M.~Maffettone, V.~V.~Gusev \emph{et al.}, ``A mobile robotic chemist,'' \emph{Nature}, vol.~583, pp.~237--241, 2020.
\bibitem{settles2009} B.~Settles, ``Active learning literature survey,'' Computer Sciences Technical Report 1648, University of Wisconsin--Madison, 2009.
\bibitem{loshchilov2019} I.~Loshchilov and F.~Hutter, ``Decoupled weight decay regularization,'' in \emph{Proc.\ ICLR}, 2019.
\end{thebibliography}

% =====================================================================
\appendices

\section{Hyperparameters}
\begin{table*}[!h]
\centering
\caption{Hyperparameters.}
\small
\begin{tabular}{lll}
\toprule
 & Stage 1 & Stages 2 / 2b \\
\midrule
Cells per neighbourhood & 64 & focal + 63 \\
Train / test & 20{,}000 / 2{,}000 patches & 46{,}423 / 1{,}811 labelled (2);\ 60{,}000 / 15{,}000 (2b) \\
Spatial split & $4\times4$ blocks & $6\times6$ blocks \\
Width / layers / heads & 128 / 4 / 4 & 128 / 4 / 4 \\
Condition & none & knockout emb.\ / expression \\
Condition dropout & -- & 0.10 / 0.15 \\
Optimiser & \multicolumn{2}{c}{AdamW, lr $3\times10^{-4}$, wd 0.01, cosine} \\
Steps / batch & 15{,}000 / 128 & 15{,}000 / 12{,}000; 128 \\
EMA & \multicolumn{2}{c}{0.999} \\
Sampler & Euler, 100 steps & Euler, 100 / 50 steps \\
Hardware & \multicolumn{2}{c}{1$\times$ NVIDIA T4, $\approx$10 min per model} \\
\bottomrule
\end{tabular}
\end{table*}

\section{Training Objective (PyTorch)}
\begin{lstlisting}
def flow_matching_loss(model, pos, typ, z, keep, n_types, pos_scale):
    # data x1 = [scaled positions, one-hot classes]; noise x0
    x1 = torch.cat([pos / pos_scale,
                    F.one_hot(typ, n_types).float()], dim=-1)
    x0 = torch.randn_like(x1)
    t = torch.rand(pos.shape[0], device=pos.device)
    tt = t[:, None, None]
    xt = (1 - tt) * x0 + tt * x1          # straight path
    v = model(xt, t, z, keep)             # keep=False -> null condition
    err = (v - (x1 - x0)) ** 2            # target velocity x1 - x0
    return err[..., :2].mean() + err[..., 2:].mean()
\end{lstlisting}

\end{document}
